\chapter{Geometry Formula Reference}
\label{app:geometry_formulas}

This appendix is a formula-level reference for the geometry used throughout
the book. It collects the exact predicates, distances, intersection
conditions, circle and sphere identities, measures, and coordinate
transformations that the algorithms rely on, each stated with its
preconditions and with pointers to the chapters that use it. Notation
follows Chapter~\ref{ch:math_foundations} and Appendix~\ref{app:math_reference}:
lowercase $p,q,r$ denote points and $u,v,w$ vectors in $\mathbb{R}^d$;
$\langle u,v\rangle$ is the dot product and $\|v\| = \sqrt{\langle v,v\rangle}$
the Euclidean norm; $u\times v$ is the scalar cross product in
$\mathbb{R}^2$ and the vector cross product in $\mathbb{R}^3$. Polygons are
oriented counterclockwise (CCW) and tetrahedra positively (right-handed)
unless stated otherwise. Standard references for this material are
\cite{deberg2008}, \cite{preparata1985}, \cite{weisstein2004}, and, for the
numerical side, \cite{higham2002}. Appendix~\ref{app:implementation_guide}
discusses when the floating-point versions of these formulas can be trusted.

\section{Vector Operations}
\label{sec:app_formulas_vectors}

\medskip\noindent\textbf{Dot product, norm, and angle.}
\index{dot product}\index{angle between vectors}
For vectors $u,v \in \mathbb{R}^d$,
\[
  \langle u, v\rangle = \sum_{i=1}^d u_i v_i,
  \qquad
  \cos\theta = \frac{\langle u,v\rangle}{\|u\|\,\|v\|},
\]
where $\theta \in [0,\pi]$ is the angle between two nonzero vectors. By
Cauchy--Schwarz (Theorem~\ref{thm:math_cauchy_schwarz}),
$|\langle u,v\rangle| \le \|u\|\,\|v\|$, so $|\cos\theta|\le 1$ and
$u \perp v$ iff $\langle u,v\rangle = 0$. The angle is obtained as
$\theta = \arccos(\langle u,v\rangle / (\|u\|\|v\|))$, computed stably by
$\operatorname{atan2}(\|u\times v\|, \langle u,v\rangle)$ in the plane.
Used in Chapter~\ref{ch:math_foundations} (Section~\ref{sec:math_inner_products})
and in every closest-point computation of Chapter~\ref{ch:distance_paths}.

\medskip\noindent\textbf{Scalar cross product in $\mathbb{R}^2$.}
\index{cross product}
For $u = (u_x, u_y)$ and $v = (v_x, v_y)$,
\[
  u\times v \;=\; u_x v_y - u_y v_x
  \;=\; \det\begin{pmatrix} u_x & v_x \\ u_y & v_y \end{pmatrix}
  \;=\; \|u\|\,\|v\|\,\sin\theta,
\]
the signed area of the parallelogram spanned by $u$ and $v$: positive when
the rotation from $u$ to $v$ is counterclockwise, zero iff $u$ and $v$ are
parallel. This is the algebraic core of the orientation test
(Section~\ref{sec:app_formulas_orientation}). Used in
Chapter~\ref{ch:geometric_predicates} and Chapter~\ref{ch:convex_hulls}.

\medskip\noindent\textbf{Vector cross product in $\mathbb{R}^3$.}
\index{cross product}
For $u, v \in \mathbb{R}^3$,
\[
  u\times v = (u_2 v_3 - u_3 v_2,\; u_3 v_1 - u_1 v_3,\; u_1 v_2 - u_2 v_1),
\]
the unique vector perpendicular to both $u$ and $v$ with
$\|u\times v\| = \|u\|\,\|v\|\,\sin\theta$ (area of the spanned
parallelogram) and direction given by the right-hand rule. Lagrange's
identity gives
\[
  \|u\times v\|^2 = \|u\|^2\|v\|^2 - \langle u,v\rangle^2,
\]
which is the numerically stable way to evaluate $\|u\times v\|$ for
near-parallel vectors. Used in Chapter~\ref{ch:polyhedral_geometry}
(normals and volumes), Chapter~\ref{ch:curves_surfaces} (surface normals),
and Chapter~\ref{ch:applications_graphics}.

\medskip\noindent\textbf{Scalar triple product.}
\index{scalar triple product}
For $u, v, w \in \mathbb{R}^3$,
\[
  \langle u,\, v\times w\rangle
  = \det\begin{pmatrix} u_x & u_y & u_z \\ v_x & v_y & v_z \\ w_x & w_y & w_z \end{pmatrix},
\]
the signed volume of the parallelepiped spanned by $u, v, w$
(Section~\ref{sec:app_math_determinants}); it is invariant under cyclic
permutations and changes sign under transpositions. Its absolute value
yields tetrahedron volumes in Section~\ref{sec:app_formulas_areas_volumes}.
Used in Chapter~\ref{ch:polyhedral_geometry}.

\medskip\noindent\textbf{Projection onto a vector.}
\index{projection}
For $v \neq \mathbf{0}$,
\[
  \operatorname{proj}_v u = \frac{\langle u,v\rangle}{\langle v,v\rangle}\,v,
  \qquad
  u = \operatorname{proj}_v u \;+\; \bigl(u - \operatorname{proj}_v u\bigr),
\]
where the residual $u - \operatorname{proj}_v u$ is orthogonal to $v$, and
$\|\operatorname{proj}_v u\| = |\langle u,v\rangle| / \|v\|$ is the length
of the orthogonal projection. Used in Chapter~\ref{ch:distance_paths}
(closest points on lines) and Chapter~\ref{ch:geometric_matching}
(projection steps in registration).

\medskip\noindent\textbf{Angle bisectors and perpendicular bisectors.}
\index{angle bisector}
For two nonzero, non-opposite vectors $u$ and $v$ the internal angle
bisector has direction
\[
  b = \frac{u}{\|u\|} + \frac{v}{\|v\|},
\]
and the external bisector has direction
$b' = u/\|u\| - v/\|v\|$, with $b \perp b'$; every vector making equal
angle with $u$ and $v$ is parallel to $b$. The locus of points equidistant
from two points $p$ and $q$ is the perpendicular bisector
\[
  \operatorname{bis}(p,q) = \bigl\{ x : \|x-p\| = \|x-q\| \bigr\}
  \;\Longleftrightarrow\;
  \langle x,\, q-p\rangle = \frac{\|q\|^2 - \|p\|^2}{2},
\]
a hyperplane perpendicular to the segment $[p,q]$ (a line in $\mathbb{R}^2$,
a plane in $\mathbb{R}^3$). For two lines given by
$\langle n_i, x\rangle = c_i$, the two angle-bisector lines satisfy
\[
  \frac{\langle n_1, x\rangle - c_1}{\|n_1\|}
  = \pm\ \frac{\langle n_2, x\rangle - c_2}{\|n_2\|},
\]
the equal-distance locus. Used in Chapter~\ref{ch:voronoi_diagrams}, where
Voronoi cells are intersections of perpendicular-bisector half-spaces and
line-site diagrams use the line bisectors (Section~\ref{sec:voronoi_segments}).

\index{vector operations}\index{perpendicular bisector}

\section{Orientation Tests}
\label{sec:app_formulas_orientation}

Orientation predicates assign to an ordered tuple of points the sign of a
determinant; they encode sidedness, collinearity, concyclicity, and
coplanarity, and they are the decisions on which most algorithms hinge.
The signs below are exactly those implemented by the robust predicates of
Chapter~\ref{ch:geometric_predicates}.

\medskip\noindent\textbf{2D orientation (sidedness).}
\index{orientation}
\index{predicate}
For points $p, q, r$ in the plane,
\[
  \operatorname{orient}(p,q,r) =
  \det\begin{pmatrix}
    1 & 1 & 1\\ p_x & q_x & r_x\\ p_y & q_y & r_y
  \end{pmatrix}
  = (q-p)\times(r-p)
  = (q_x-p_x)(r_y-p_y) - (q_y-p_y)(r_x-p_x).
\]
The value is positive iff $r$ lies strictly to the left of the directed
line $p\to q$ (the turn $p, q, r$ is CCW), negative iff to the right, and
zero iff $p, q, r$ are collinear. It equals twice the signed area of the
triangle $pqr$ (Section~\ref{sec:math_signed_area_volume}). Used in
Chapter~\ref{ch:geometric_predicates} (Section~\ref{sec:orientation_predicate}),
Chapter~\ref{ch:convex_hulls} (Sections~\ref{sec:graham_scan} and
\ref{sec:monotone_chain}), Chapter~\ref{ch:segment_intersection}, and
Chapter~\ref{ch:point_location}.

\medskip\noindent\textbf{3D orientation (coplanarity).}
\index{orientation}
For points $p, q, r, s$ in $\mathbb{R}^3$,
\[
  \operatorname{orient}(p,q,r,s) =
  \det\begin{pmatrix}
    1 & 1 & 1 & 1\\
    p_x & q_x & r_x & s_x\\
    p_y & q_y & r_y & s_y\\
    p_z & q_z & r_z & s_z
  \end{pmatrix}
  = \bigl\langle (q-p)\times(r-p),\, s-p\bigr\rangle,
\]
which is six times the signed volume of the tetrahedron $pqrs$. The value
is positive iff $s$ lies on the positive side of the oriented plane through
$p, q, r$ and zero iff the four points are coplanar. Used in
Chapter~\ref{ch:geometric_predicates} (Section~\ref{sec:orientation3d_predicate}),
Chapter~\ref{ch:convex_polytopes}, and Chapter~\ref{ch:mesh_generation}.

\medskip\noindent\textbf{In-circle test.}
\index{in-circle test}
For $a, b, c, d$ with $a, b, c$ counterclockwise,
\[
  \operatorname{incircle}(a,b,c,d) =
  \det\begin{pmatrix}
    a_x & a_y & a_x^2 + a_y^2 & 1\\
    b_x & b_y & b_x^2 + b_y^2 & 1\\
    c_x & c_y & c_x^2 + c_y^2 & 1\\
    d_x & d_y & d_x^2 + d_y^2 & 1
  \end{pmatrix},
\]
which is positive iff $d$ lies strictly inside the circle through
$a, b, c$. It is the sign of the position of the lifted point
$\hat d = (d_x, d_y, d_x^2 + d_y^2)$ with respect to the plane through the
lifted points $\hat a, \hat b, \hat c$ on the paraboloid
$z = x^2 + y^2$ (Section~\ref{sec:delaunay_lifting}). Used in
Chapter~\ref{ch:delaunay_triangulations} (Section~\ref{sec:incircle_predicate},
Definition~\ref{def:incircle}) and Chapter~\ref{ch:geometric_predicates}.

\medskip\noindent\textbf{In-sphere test.}
\index{in-sphere test}
The three-dimensional analogue for $a, b, c, d, e$,
\[
  \operatorname{insphere}(a,b,c,d,e) =
  \det\begin{pmatrix}
    a_x & a_y & a_z & a_x^2+a_y^2+a_z^2 & 1\\
    b_x & b_y & b_z & b_x^2+b_y^2+b_z^2 & 1\\
    c_x & c_y & c_z & c_x^2+c_y^2+c_z^2 & 1\\
    d_x & d_y & d_z & d_x^2+d_y^2+d_z^2 & 1\\
    e_x & e_y & e_z & e_x^2+e_y^2+e_z^2 & 1
  \end{pmatrix},
\]
is positive iff $e$ lies strictly inside the sphere through $a, b, c, d$,
for a positively oriented quadruple $a, b, c, d$; it is the in-circle
predicate lifted one dimension. Used in Chapter~\ref{ch:geometric_predicates}
(Section~\ref{sec:three_dimensional_predicates}) and in 3D Delaunay
constructions of Chapter~\ref{ch:delaunay_triangulations}.

A zero determinant always marks a degenerate configuration: collinear,
coplanar, concyclic, or cospherical points. Algorithms therefore assume
general position and handle degeneracies separately
(Chapter~\ref{ch:introduction}, Section~\ref{sec:intro_degeneracies}).

\subsection{Robustness Constants and Error Bounds}
\label{sec:appform_robustness}
\index{machine epsilon}\index{unit roundoff}\index{floating-point}

In binary floating point with a $p$-bit significand, arithmetic is exact
only up to a relative error bounded by the \emph{unit roundoff}
$u = 2^{-p}$: for $\odot \in \{+, -, \times, \div\}$,
\[
  \operatorname{fl}(a \odot b) = (a \odot b)\,(1 + \delta),
  \qquad |\delta| \le u,
\]
and the \emph{machine epsilon} is $\varepsilon = 2u = 2^{1-p}$. The standard
precisions are collected below (\cite{higham2002}, Chapter 2):
\begin{center}
\begin{tabular}{llll}
\toprule
format & significand $p$ & unit roundoff $u$ & machine epsilon $\varepsilon$ \\
\midrule
binary32 (single)  & $24$  & $2^{-24} \approx 5.96\times10^{-8}$  & $2^{-23} \approx 1.19\times10^{-7}$ \\
binary64 (double)  & $53$  & $2^{-53} \approx 1.11\times10^{-16}$ & $2^{-52} \approx 2.22\times10^{-16}$ \\
extended (80-bit)  & $64$  & $2^{-64} \approx 5.42\times10^{-20}$ & $2^{-63} \approx 1.08\times10^{-19}$ \\
binary128 (quad)   & $113$ & $2^{-113} \approx 9.63\times10^{-35}$ & $2^{-112} \approx 1.93\times10^{-34}$ \\
\bottomrule
\end{tabular}
\end{center}

A useful shorthand is the factor
\[
  \gamma_n = \frac{n\,u}{1 - n\,u},
\]
which bounds the accumulated relative error of a computation of $n$
floating-point operations. For example, the computed dot product of two
$n$-vectors satisfies
\[
  \bigl|\operatorname{fl}\bigl(\textstyle\sum_{i=1}^n x_i y_i\bigr)
  - \langle x, y\rangle\bigr|
  \le \gamma_n \sum_{i=1}^n |x_i y_i|
\]
(\cite{higham2002}, Chapter 3), so cancellation of nearly equal
large-magnitude terms is the danger.

For the 2D orientation test with coordinates bounded by $B$
($|p_x|,|p_y|,\dots \le B$), the seven operations of the standard formula
yield the first-order absolute error bound
\[
  \bigl|\widehat{\operatorname{orient}}(p,q,r) - \operatorname{orient}(p,q,r)\bigr|
  \le 28\,u\,B^2 + O(u^2 B^2);
\]
the sign is therefore reliable whenever $|\operatorname{orient}|$ exceeds
the bound, which is the principle behind the error filters and adaptive
exact evaluation of Section~\ref{sec:floating_point_filters} and
Section~\ref{sec:robust_predicates}; the tight constants are worked out by
\cite{shewchuk1997}. In Appendix~\ref{app:implementation_guide} the same
constants appear as the predicates' epsilon thresholds
(Sections~\ref{sec:appimpl_machine_eps} and \ref{sec:appimpl_adaptive}).

For a determinant, the condition number under componentwise perturbations
is the ratio of the sum of the absolute values of the permutation products
to the determinant itself: writing $a_{ij}$ for the entries of an
$n\times n$ matrix $A$ and $\widehat a_{ij} = a_{ij}(1+\varepsilon_{ij})$
with $|\varepsilon_{ij}| \le u$,
\[
  |\det\widehat A - \det A|
  \le u \sum_{\sigma \in S_n}\prod_{i=1}^n |a_{i,\sigma(i)}| + O(u^2),
\]
so the relative condition number of the determinant is
\[
  \kappa = \frac{\sum_{\sigma \in S_n}\prod_i |a_{i,\sigma(i)}|}{|\det A|},
\]
which is enormous precisely when the rows or columns are nearly dependent,
i.e., when the geometric configuration is nearly degenerate
(\cite{higham2002}, Section 12.3). This is why near-zero determinants must
not be evaluated in floating point without a filter; exact arithmetic
(Section~\ref{sec:exact_arithmetic}) removes the problem entirely.

\index{orientation tests}\index{condition number}

\section{Distance Formulas}
\label{sec:app_formulas_distances}

\medskip\noindent\textbf{Point to line.}
\index{point-line distance}
For a point $q$ and a line $\ell = \{ p + t v : t \in \mathbb{R} \}$ with
$v \neq \mathbf{0}$,
\[
  \operatorname{dist}(q, \ell) = \frac{\|(q - p)\times v\|}{\|v\|},
\]
using the scalar cross product in $\mathbb{R}^2$ and the vector cross
product in $\mathbb{R}^3$. Equivalently, for a line in normal form
$\langle n, x\rangle = c$, the signed distance is
$\delta(q) = (\langle n, q\rangle - c)/\|n\|$ and
$\operatorname{dist}(q, \ell) = |\delta(q)|$ (Section~\ref{sec:math_lines_hyperplanes}).
Used in Chapter~\ref{ch:distance_paths} and in the walking point-location
of Chapter~\ref{ch:point_location}.

\medskip\noindent\textbf{Point to plane.}
\index{point-plane distance}
For a point $q$ and a plane $H = \{ x : \langle n, x\rangle = c \}$ with
$n \neq \mathbf{0}$,
\[
  \operatorname{dist}(q, H) = \frac{|\langle n, q\rangle - c|}{\|n\|},
\]
with signed version $\delta(q) = (\langle n, q\rangle - c)/\|n\|$
positive on the side pointed to by $n$. The same formula (with
$n = (a,b)$, $c = -d$) covers a line in $\mathbb{R}^2$. Used in
Chapter~\ref{ch:distance_paths} and Chapter~\ref{ch:convex_polytopes}.

\medskip\noindent\textbf{Point to segment.}
\index{point-segment distance}
For a point $q$ and the segment $[p,\, p+v]$, $v \neq \mathbf{0}$, the
closest point is $p + t^* v$ with the projection parameter clamped to the
segment:
\[
  t^* = \min\Bigl(1,\ \max\Bigl(0,\ \frac{\langle q-p, v\rangle}{\langle v, v\rangle}\Bigr)\Bigr),
  \qquad
  \operatorname{dist}\bigl(q,\,[p,\, p+v]\bigr) = \|q - p - t^* v\|.
\]
Used in Chapter~\ref{ch:distance_paths}, Chapter~\ref{ch:proximity}, and
Chapter~\ref{ch:motion_planning} (distance to a robot link).

\medskip\noindent\textbf{Between two lines in $\mathbb{R}^3$.}
\index{line-line distance}
For skew lines $p + t v$ and $q + s w$ with $v\times w \neq \mathbf{0}$,
the distance between the lines is
\[
  \operatorname{dist}(\ell_1, \ell_2)
  = \frac{\bigl|\langle q - p,\, v\times w\rangle\bigr|}{\|v\times w\|},
\]
the absolute scalar triple product divided by the area of the direction
parallelogram; the closest points are $p + t^* v$ and $q + s^* w$ where
$(t^*, s^*)$ solves the $2\times 2$ system of
Section~\ref{sec:app_formulas_intersections}. For parallel lines
($v\times w = \mathbf{0}$) the distance reduces to the point-to-line
distance from any point of one line to the other. Used in
Chapter~\ref{ch:distance_paths} and Chapter~\ref{ch:convex_intersection}.

\medskip\noindent\textbf{Between two segments (2D and 3D).}
\index{segment-segment distance}
The distance between segments $[p, p+v]$ and $[q, q+w]$ is the minimum
\[
  \min_{s,t \in [0,1]} \bigl\| p + t\,v - q - s\,w \bigr\|^2,
\]
a convex quadratic in $(t, s)$. Its interior critical point solves
\[
  \begin{pmatrix}
    \langle v,v\rangle & -\langle v,w\rangle\\
    -\langle v,w\rangle & \langle w,w\rangle
  \end{pmatrix}
  \begin{pmatrix} t \\ s \end{pmatrix}
  = \begin{pmatrix} \langle v,\, q-p\rangle \\ \langle w,\, q-p\rangle \end{pmatrix},
\]
and the minimum over the box $[0,1]^2$ is either this point (when it lies
inside) or is attained on the boundary, which is a point-segment distance;
the closest points are returned along with the distance. In $\mathbb{R}^2$
the two segments intersect iff this minimum is zero, which is decided
directly by the orientation products of Section~\ref{sec:app_formulas_intersections}.
Used in Chapter~\ref{ch:distance_paths} (closest-pair-adjacent queries)
and Chapter~\ref{ch:motion_planning} (clearance of link chains).

\medskip\noindent\textbf{Between parallel planes.}
For parallel planes $\langle n, x\rangle = c_1$ and
$\langle n, x\rangle = c_2$ (same normal, up to scaling),
\[
  \operatorname{dist}(H_1, H_2) = \frac{|c_1 - c_2|}{\|n\|}.
\]
Used in Chapter~\ref{ch:convex_polytopes} and
Chapter~\ref{ch:applications_graphics} (slab tests).

\index{distances}\index{closest point}

\section{Intersection Formulas}
\label{sec:app_formulas_intersections}

\medskip\noindent\textbf{Two lines in $\mathbb{R}^2$.}
\index{line-line intersection}
For lines $p + t v$ and $q + s w$ with $v \times w \neq 0$,
\[
  t = \frac{(q - p)\times w}{v \times w},
  \qquad
  s = \frac{(q - p)\times v}{v \times w},
\]
and the intersection point is $p + t v = q + s w$. If
$v \times w = 0$ the lines are parallel: coincident (no unique
intersection) or disjoint. Used in Chapter~\ref{ch:segment_intersection},
Chapter~\ref{ch:geometric_duality}, and Chapter~\ref{ch:applications_graphics}.

\medskip\noindent\textbf{Two segments in $\mathbb{R}^2$.}
\index{segment intersection}
The closed segments $[a, b]$ and $[c, d]$ intersect iff the endpoints of
each straddle the line of the other:
\[
  \operatorname{orient}(a,b,c)\cdot\operatorname{orient}(a,b,d) \le 0
  \qquad\text{and}\qquad
  \operatorname{orient}(c,d,a)\cdot\operatorname{orient}(c,d,b) \le 0,
\]
where equality holds exactly for boundary contact (an endpoint on the other
segment, or overlapping collinear segments, which are then reduced to a
1D overlap test). Used in Chapter~\ref{ch:segment_intersection}
(Section~\ref{sec:segint_pairwise}).

\medskip\noindent\textbf{Line and plane in $\mathbb{R}^3$.}
\index{line-plane intersection}
For the line $p + t v$ and the plane $\langle n, x\rangle = c$,
\[
  t = \frac{c - \langle n, p\rangle}{\langle n, v\rangle},
  \qquad\text{provided } \langle n, v\rangle \neq 0,
\]
and the intersection point is $p + t v$; if $\langle n, v\rangle = 0$ the
line is parallel to the plane (contained or disjoint). Used in
Chapter~\ref{ch:applications_graphics} (Section~\ref{sec:graphics_ray_intersection})
and Chapter~\ref{ch:convex_polytopes}.

\medskip\noindent\textbf{Two planes in $\mathbb{R}^3$.}
\index{plane-plane intersection}
The planes $\langle n_1, x\rangle = c_1$ and $\langle n_2, x\rangle = c_2$
with $n_1\times n_2 \neq \mathbf{0}$ meet in a line with direction
\[
  d = n_1 \times n_2
\]
through any point $p_0$ satisfying the two equations; a convenient choice
solves the $2\times 3$ system
\[
  \begin{pmatrix} n_1^\top \\ n_2^\top \end{pmatrix} p_0
  = \begin{pmatrix} c_1 \\ c_2 \end{pmatrix}
\]
for the component of $p_0$ along the direction of largest pivot
(Section~\ref{sec:app_math_linear_algebra}). If $n_1\times n_2 = \mathbf{0}$
the planes are parallel. Used in Chapter~\ref{ch:convex_intersection}
(separating planes) and Chapter~\ref{ch:convex_polytopes} (facet-plane
interactions).

\medskip\noindent\textbf{Ray and triangle (M\"oller--Trumbore).}
\index{ray-triangle intersection}
Let the triangle have vertices $v_0, v_1, v_2$, ray origin $O$, and unit
direction $d$. Writing the hit point in barycentric coordinates of
Section~\ref{sec:app_formulas_barycentric}, the ray intersects the triangle
where
\[
  E_1\,u + E_2\,v + d\,t = T,
  \qquad
  E_1 = v_1 - v_0,\quad E_2 = v_2 - v_0,\quad T = O - v_0,
\]
which Cramer's rule solves as
\[
  u = \frac{\langle T,\, d\times E_2\rangle}{\langle E_1,\, d\times E_2\rangle},
  \qquad
  v = \frac{\langle d,\, T\times E_1\rangle}{\langle E_1,\, d\times E_2\rangle},
  \qquad
  t = \frac{\langle E_2,\, T\times E_1\rangle}{\langle E_1,\, d\times E_2\rangle}.
\]
The ray hits iff $t \ge 0$, $u \ge 0$, $v \ge 0$, and $u + v \le 1$, with
$\langle E_1, d\times E_2\rangle = 0$ the parallel (miss) case. Used in
Chapter~\ref{ch:applications_graphics} (Section~\ref{sec:graphics_ray_intersection})
and in the Monte Carlo volume estimates of Chapter~\ref{ch:sampling}.

\subsection{Arrangements and Counting Formulas}
\label{sec:appform_counting}
\index{arrangement}\index{Euler formula}

\medskip\noindent\textbf{Arrangement of lines.}
For an arrangement of $n$ lines in general position (no two parallel, no
three concurrent), the numbers of cells are
\[
  V = \binom{n}{2} = \frac{n(n-1)}{2},
  \qquad
  E = n^2,
  \qquad
  F = \frac{n(n+1)}{2} + 1
\]
vertices, edges, and faces (including the unbounded face); here the graph
relation reads $V - E + F = 1$ because the unbounded rays have a single
endpoint, and treating the point at infinity restores $V - E + F = 2$
(\cite{deberg2008}, Chapter 8). Used in Chapter~\ref{ch:segment_intersection}
(Section~\ref{sec:line_arrangements}), Chapter~\ref{ch:envelopes}, and
Chapter~\ref{ch:geometric_duality}.

\begin{theorem}[Zone Theorem]\label{thm:appform_zone}
  In an arrangement of $n$ lines the zone of any line has $O(n)$ edges;
  in an arrangement of $n$ line segments the zone of a segment has
  complexity $O(n\,\alpha(n))$, where $\alpha$ is the inverse
  Ackermann function (\cite{deberg2008}, Chapter 8).
\end{theorem}
The zone theorem drives the complexity of the sweep of
Chapter~\ref{ch:segment_intersection} and the randomized incremental
constructions of Chapter~\ref{ch:randomized_algorithms}.

\medskip\noindent\textbf{Planar graphs.}
\index{planar graph}
For a connected planar map with $V$ vertices, $E$ edges, and $F$ faces
(including the unbounded face), Euler's formula holds:
\[
  V - E + F = 2
\]
(\cite{euler1758}); equivalently $E \le 3V - 6$ and $F \le 2V - 4$ for a
triangulation of $V \ge 3$ vertices. Used in Chapter~\ref{ch:computational_topology}
(Section~\ref{sec:euler_characteristic}), Chapter~\ref{ch:half_edge}
(Section~\ref{sec:halfedge_complexity}), and Chapter~\ref{ch:mesh_analysis}.

\medskip\noindent\textbf{Triangulation counts.}
A triangulation of a simple polygon with $n$ vertices has
\[
  T = n - 2 \text{ triangles}, \qquad
  E = 2n - 3 \text{ edges (incl.\ boundary)}, \qquad
  n - 3 \text{ diagonals}.
\]
A triangulation of a point set of $n$ points, $h$ of them on the convex
hull, has
\[
  T = 2n - 2 - h \text{ triangles}, \qquad
  E = 3n - 3 - h \text{ edges}.
\]
Used in Chapter~\ref{ch:triangulations}, Chapter~\ref{ch:polygon_partitioning},
Chapter~\ref{ch:art_gallery}, and Chapter~\ref{ch:delaunay_triangulations}.

\medskip\noindent\textbf{Series used in complexity analysis.}
\index{series}\index{harmonic number}
The following sums appear throughout the randomized and expected-time
analyses of Chapters~\ref{ch:randomized_algorithms},
\ref{ch:delaunay_triangulations}, and \ref{ch:high_dimensional}
(\cite{cormen2009}, Appendix A):
\begin{center}
\begin{tabular}{ll}
\toprule
$\displaystyle\sum_{i=1}^{n} i = \frac{n(n+1)}{2}$ &
  arithmetic sum \\
$\displaystyle\sum_{i=1}^{n} i^2 = \frac{n(n+1)(2n+1)}{6}$ &
  squared arithmetic sum \\
$\displaystyle\sum_{i=0}^{n} 2^i = 2^{n+1} - 1$ &
  doubling sum \\
$\displaystyle\sum_{i=0}^{n} r^i = \frac{1 - r^{n+1}}{1 - r}$, $r \neq 1$ &
  geometric series \\
$\displaystyle\sum_{i=1}^{n} \frac{1}{i} = H_n = \ln n + \gamma + O\!\Bigl(\frac{1}{n}\Bigr)$ &
  harmonic number, $\gamma \approx 0.5772$ \\
$\displaystyle\sum_{i=1}^{\infty} \frac{1}{i^2} = \frac{\pi^2}{6}$ &
  Euler's sum \\
$\displaystyle\sum_{k=0}^{n} \binom{n}{k} = 2^n$, 
  $\displaystyle\sum_{k=0}^{n} \binom{n}{k} x^k = (1+x)^n$ &
  binomial theorem \\
$\displaystyle\sum_{i=k}^{n} \binom{i}{k} = \binom{n+1}{k+1}$ &
  hockey-stick identity \\
\bottomrule
\end{tabular}
\end{center}

\index{intersection formulas}\index{counting}

\section{Circle Predicates}
\label{sec:app_formulas_circles}

\medskip\noindent\textbf{Circumcircle of three points.}
\index{circumcircle}
For three non-collinear points $a, b, c$, the circumcenter $o$ is the
unique point equidistant from all three, solving
\[
  2 \begin{pmatrix} (b - a)^\top \\ (c - a)^\top \end{pmatrix} o
  = \begin{pmatrix} \|b\|^2 - \|a\|^2 \\ \|c\|^2 - \|a\|^2 \end{pmatrix},
\]
and the circumradius is $r = \|o - a\|$. Expanding Cramer's rule gives the
explicit coordinates
\[
  D = 2\,\bigl(x_a(y_b - y_c) + x_b(y_c - y_a) + x_c(y_a - y_b)\bigr) \neq 0,
\]
\[
  o_x = \frac{(x_a^2 + y_a^2)(y_b - y_c) + (x_b^2 + y_b^2)(y_c - y_a) + (x_c^2 + y_c^2)(y_a - y_b)}{D},
\]
\[
  o_y = \frac{(x_a^2 + y_a^2)(x_c - x_b) + (x_b^2 + y_b^2)(x_a - x_c) + (x_c^2 + y_c^2)(x_b - x_a)}{D},
\]
where $D = 2\,\operatorname{orient}(a,b,c)$ equals four times the signed
area of the triangle; small $D$ means a skinny triangle and an
ill-conditioned circumcenter
(\cite{weisstein2004}). Used in Chapter~\ref{ch:delaunay_triangulations}
(Section~\ref{sec:delaunay_empty_circumcircle}), Chapter~\ref{ch:voronoi_diagrams}
(circumcenters are Voronoi vertices), and Chapter~\ref{ch:sampling}
(Section~\ref{sec:largest_empty_circle}).

\medskip\noindent\textbf{Circumsphere of four points.}
\index{circumsphere}
For four non-coplanar points $p_0, p_1, p_2, p_3$, the circumcenter $o$
solves the $3\times 3$ system
\[
  2 \begin{pmatrix} (p_1 - p_0)^\top \\ (p_2 - p_0)^\top \\ (p_3 - p_0)^\top \end{pmatrix} o
  = \begin{pmatrix} \|p_1\|^2 - \|p_0\|^2 \\
      \|p_2\|^2 - \|p_0\|^2 \\
      \|p_3\|^2 - \|p_0\|^2 \end{pmatrix},
\]
with radius $r = \|o - p_0\|$. Used in Chapter~\ref{ch:delaunay_triangulations}
(empty-circumsphere property of 3D Delaunay) and
Chapter~\ref{ch:mesh_generation}. Whether a point lies inside the sphere is
decided exactly by the in-sphere determinant of
Section~\ref{sec:app_formulas_orientation}.

\medskip\noindent\textbf{Circle--circle intersection.}
\index{circle-circle intersection}
Two circles with centers $c_1, c_2$ and radii $r_1, r_2$, at distance
$d = \|c_2 - c_1\|$, intersect iff
\[
  |r_1 - r_2| \le d \le r_1 + r_2,
\]
with equality in either bound a tangency. For $|r_1 - r_2| < d < r_1 + r_2$
there are two intersection points $p_\pm$ given by
\[
  a = \frac{d^2 + r_1^2 - r_2^2}{2d},
  \qquad
  h = \sqrt{r_1^2 - a^2},
  \qquad
  p_\pm = c_1 + \frac{a}{d}\,(c_2 - c_1) \pm \frac{h}{d}\,R_{90}(c_2 - c_1),
\]
where $R_{90}(u) = (-u_y, u_x)$ is the CCW rotation by $90^\circ$. Used in
Chapter~\ref{ch:voronoi_diagrams} (Voronoi vertices as circle--circle
intersections) and Chapter~\ref{ch:proximity}.

\medskip\noindent\textbf{Circle--line intersection.}
\index{circle-line intersection}
For a circle with center $c$ and radius $r$, and a line through $p$ with
unit direction $v$, the projection parameter of the closest point to $c$
is $t = \langle c - p, v\rangle$; the intersections, if any, are
\[
  p_\pm = p + \bigl(t \pm \sqrt{\Delta}\bigr) v,
  \qquad
  \Delta = r^2 - \|c - p\|^2 + t^2,
\]
with two distinct intersections for $\Delta > 0$, one tangent point for
$\Delta = 0$, and none for $\Delta < 0$. Restricting to $t \ge 0$ yields the
ray--circle (ray--sphere in $\mathbb{R}^3$) intersection. Used in
Chapter~\ref{ch:applications_graphics} and Chapter~\ref{ch:packing_algorithms}.

\medskip\noindent\textbf{Power of a point.}
\index{power of a point}
For a point $p$ and a circle with center $c$ and radius $r$, the power is
\[
  \operatorname{Pow}(p) = \|p - c\|^2 - r^2,
\]
positive outside, negative inside, and zero on the circle. For any line
through $p$ meeting the circle at $X$ and $Y$,
\[
  \operatorname{Pow}(p) = \langle p - X,\, p - Y\rangle,
\]
the product of the signed directed distances, independent of the choice of
line. Used in Chapter~\ref{ch:voronoi_diagrams} (power diagrams and
weighted Voronoi cells) and Chapter~\ref{ch:geometric_duality}.

\medskip\noindent\textbf{Radical axis.}
\index{radical axis}
The set of points with equal power with respect to two circles
$(c_1, r_1)$ and $(c_2, r_2)$ is the line
\[
  \|x - c_1\|^2 - r_1^2 = \|x - c_2\|^2 - r_2^2
  \;\Longleftrightarrow\;
  2\,\langle x,\, c_2 - c_1\rangle = \|c_2\|^2 - \|c_1\|^2 + r_1^2 - r_2^2,
\]
perpendicular to the line of centers; it is the common chord when the two
circles intersect. Used in Chapter~\ref{ch:voronoi_diagrams} (edges of
power diagrams) and Chapter~\ref{ch:geometric_duality}.

\index{circle predicates}\index{circumcenter}

\section{Areas and Volumes}
\label{sec:app_formulas_areas_volumes}

\medskip\noindent\textbf{Triangle area in $\mathbb{R}^2$.}
\index{triangle area}
For a triangle with vertices $p, q, r$,
\[
  A = \frac{1}{2}\,\bigl|(q - p)\times(r - p)\bigr|
  = \frac{1}{2}\,\bigl|\operatorname{orient}(p,q,r)\bigr|,
\]
or, in terms of the side lengths $a, b, c$, Heron's formula
\[
  A = \sqrt{s(s-a)(s-b)(s-c)},
  \qquad
  s = \frac{a + b + c}{2}
\]
(\cite{weisstein2004}). Used in Chapter~\ref{ch:polygon_ops} and
Chapter~\ref{ch:mesh_analysis} (element area and quality).

\medskip\noindent\textbf{Polygon area (shoelace formula).}
\index{shoelace formula}\index{polygon area}
For a simple polygon $p_1, \dots, p_n$ (indices modulo $n$), the signed
area is
\[
  A = \frac{1}{2}\sum_{i=1}^{n} (x_i y_{i+1} - x_{i+1} y_i),
\]
positive for a CCW boundary and equal in absolute value to the usual area
(Section~\ref{sec:math_signed_area_volume}). The formula extends to
self-intersecting polygons as the signed area enclosed by the boundary.
Used in Chapter~\ref{ch:polygon_ops} (area, orientation repair) and
Chapter~\ref{ch:surface_geometry}.

\medskip\noindent\textbf{Polygon centroid.}
\index{centroid}
The centroid of a triangle is the arithmetic mean of its vertices. The
area-weighted centroid of a polygon is the average of the triangle
centroids with the signed areas as weights:
\[
  G = \frac{1}{6A}\sum_{i=1}^{n} (x_i y_{i+1} - x_{i+1} y_i)\,(p_i + p_{i+1}),
\]
with $A$ the signed area above; note that the arithmetic mean of the
vertices equals $G$ only for triangles. Used in Chapter~\ref{ch:polygon_ops}
(area, moment, and center-of-mass computations).

\medskip\noindent\textbf{Convex hull perimeter and isoperimetric inequality.}
\index{convex hull}\index{isoperimetric inequality}
The perimeter of a convex polygon with vertices $p_1, \dots, p_n$ in
boundary order is
\[
  \operatorname{perim}(P) = \sum_{i=1}^{n}\|p_{i+1} - p_i\|,
\]
the length of the boundary of the convex hull of a point set
(Chapter~\ref{ch:convex_hulls}). For any closed curve enclosing area $A$,
\[
  L^2 \ge 4\pi A,
  \qquad
  L \ge 2\sqrt{\pi A},
\]
the isoperimetric inequality, with equality only for a circle; it governs
the packing-density arguments of Chapter~\ref{ch:packing_algorithms}.

\medskip\noindent\textbf{Tetrahedron volume.}
\index{tetrahedron volume}
For a tetrahedron with vertices $p, q, r, s$,
\[
  V = \frac{1}{6}\,\bigl|\bigl\langle (q-p)\times(r-p),\, s-p\bigr\rangle\bigr|
  = \frac{1}{6}\,\bigl|\operatorname{orient}(p,q,r,s)\bigr|,
\]
one sixth of the scalar triple product (Section~\ref{sec:app_formulas_vectors}).
Used in Chapter~\ref{ch:polyhedral_geometry}, Chapter~\ref{ch:mesh_generation}
(element volumes), and Chapter~\ref{ch:mesh_analysis}.

\medskip\noindent\textbf{Polyhedron volume (divergence theorem).}
\index{polyhedron volume}\index{divergence theorem}
For a closed surface mesh with consistently oriented triangular faces
$(v_1, v_2, v_3)$ (CCW seen from outside), the enclosed volume is
\[
  V = \frac{1}{6}\sum_{f} \bigl\langle v_1,\ v_2\times v_3\bigr\rangle,
\]
the sum of the signed tetrahedra formed by each face with the origin
(any fixed apex works); this is the discrete divergence theorem
(Section~\ref{sec:app_math_vector_calculus}). Used in
Chapter~\ref{ch:polyhedral_geometry} (volume, center of mass, moments) and
Chapter~\ref{ch:sampling} (validation of Monte Carlo estimates).

\medskip\noindent\textbf{Surface area of a mesh.}
\index{surface area}
The surface area of a triangle mesh is the sum of the triangle areas
\[
  S = \sum_{f} \frac{1}{2}\,\bigl\|(v_2 - v_1)\times(v_3 - v_1)\bigr\|,
\]
and the area of a triangulated patch is orientation-independent. Used in
Chapter~\ref{ch:surface_geometry} and Chapter~\ref{ch:mesh_analysis}.

\medskip\noindent\textbf{Centroid (center of mass) of a polyhedron.}
\index{centroid}
The centroid of a tetrahedron is the arithmetic mean of its four vertices.
For a closed triangle mesh with consistently oriented faces, the centroid
is the signed-volume-weighted average of the tetrahedron centroids:
\[
  G = \frac{1}{4V}\sum_{f} V_f\,(v_1 + v_2 + v_3),
  \qquad
  V_f = \frac{1}{6}\,\langle v_1, v_2\times v_3\rangle,
\]
with $V = \sum_f V_f$. Used in Chapter~\ref{ch:polyhedral_geometry} and
Chapter~\ref{ch:applications_robotics} (balance and inertial properties).

\medskip\noindent\textbf{Simplex volume in $d$ dimensions.}
\index{simplex volume}
The $d$-simplex $\Delta^d = \operatorname{conv}\{p_0, \dots, p_d\}$ has
$d$-dimensional volume
\[
  \operatorname{Vol}(\Delta^d)
  = \frac{1}{d!}\,
  \biggl|\det\begin{pmatrix} p_1 - p_0 \\ \vdots \\ p_d - p_0 \end{pmatrix}\biggr|,
\]
generalizing the $1/2$ and $1/6$ factors above; the standard simplex
$\operatorname{conv}\{0, e_1, \dots, e_d\}$ has volume $1/d!$. Used in
Chapter~\ref{ch:high_dimensional} (high-dimensional volume) and
Chapter~\ref{ch:sampling} (volume estimation).

\medskip\noindent\textbf{Volumes of balls and spheres.}
\index{ball volume}
The volume of the Euclidean ball of radius $r$ in $\mathbb{R}^d$ and the
surface volume of its boundary sphere are
\[
  V_d(r) = \frac{\pi^{d/2}}{\Gamma(d/2 + 1)}\,r^d,
  \qquad
  S_d(r) = \frac{d}{r}\,V_d(r),
\]
so in $\mathbb{R}^3$, $V_3(r) = \frac{4}{3}\pi r^3$ and
$S_3(r) = 4\pi r^2$. The concentration of volume near the boundary in high
dimension drives the phenomena of Chapter~\ref{ch:high_dimensional}.

\index{areas and volumes}\index{volume}

\section{Coordinate Transformations}
\label{sec:app_formulas_transformations}

\medskip\noindent\textbf{Affine maps as matrices.}
\index{affine transformation}
An affine map $x \mapsto Mx + t$ (invertible $d\times d$ matrix $M$,
translation $t$) preserves collinearity, parallelism, ratios along lines,
and affine combinations, but not lengths or angles in general
(Section~\ref{sec:math_transformations}). In homogeneous coordinates it is
the linear map (Section~\ref{sec:math_homogeneous})
\[
  \begin{pmatrix} \tilde x \\ w \end{pmatrix}
  \mapsto
  \begin{pmatrix} M & t \\ \mathbf{0}^\top & 1 \end{pmatrix}
  \begin{pmatrix} \tilde x \\ w \end{pmatrix},
  \qquad
  x = \tilde x / w,
\]
so composition of affine maps is ordinary matrix multiplication. Used in
Chapter~\ref{ch:applications_graphics} (model, view, and projection
matrices), Chapter~\ref{ch:applications_robotics}, and
Chapter~\ref{ch:geometric_matching}.

\medskip\noindent\textbf{Translation and scaling.}
\index{translation}\index{scaling}
Translation by $t$ is $(I, t)$; uniform scaling by $s$ is $(sI, \mathbf{0})$;
anisotropic scaling about the origin is
$\operatorname{diag}(s_1, \dots, s_d)$. Scaling or rotation about a point
$p$ is the conjugation
\[
  T(p)\,M\,T(-p),
\]
with $T(\cdot)$ the translation matrix. Used in Chapter~\ref{ch:applications_graphics}
and Chapter~\ref{ch:geometric_matching} (shape alignment).

\medskip\noindent\textbf{Rotation in $\mathbb{R}^2$.}
\index{rotation}
The CCW rotation by angle $\theta$ is
\[
  R(\theta) = \begin{pmatrix}
    \cos\theta & -\sin\theta\\
    \sin\theta & \cos\theta
  \end{pmatrix},
\]
with $R(\theta_1)R(\theta_2) = R(\theta_1 + \theta_2)$,
$R(-\theta) = R(\theta)^\top = R(\theta)^{-1}$, and $\det R = 1$. Used in
Chapter~\ref{ch:geometric_matching} (Section~\ref{sec:matching_rigid}) and
Chapter~\ref{ch:applications_graphics}.

\medskip\noindent\textbf{Rotation in $\mathbb{R}^3$.}
\index{rotation}
The rotations about the coordinate axes are
\[
  R_x(\theta) = \begin{pmatrix}
    1 & 0 & 0\\
    0 & \cos\theta & -\sin\theta\\
    0 & \sin\theta & \cos\theta
  \end{pmatrix}, \quad
  R_y(\theta) = \begin{pmatrix}
    \cos\theta & 0 & \sin\theta\\
    0 & 1 & 0\\
    -\sin\theta & 0 & \cos\theta
  \end{pmatrix}, \quad
  R_z(\theta) = \begin{pmatrix}
    \cos\theta & -\sin\theta & 0\\
    \sin\theta & \cos\theta & 0\\
    0 & 0 & 1
  \end{pmatrix}.
\]
Rotation by angle $\theta$ about a unit axis $k$ is given by Rodrigues'
formula
\[
  R = I + \sin\theta\,[k]_\times + (1 - \cos\theta)\,[k]_\times^2,
  \qquad
  [k]_\times = \begin{pmatrix}
    0 & -k_z & k_y\\
    k_z & 0 & -k_x\\
    -k_y & k_x & 0
  \end{pmatrix},
\]
where $[k]_\times$ is the skew-symmetric matrix of the cross product
$[k]_\times v = k\times v$. Used in Chapter~\ref{ch:applications_robotics}
(robot kinematics) and Chapter~\ref{ch:geometric_matching} (rigid
registration).

\medskip\noindent\textbf{Change of basis.}
\index{change of basis}
If the columns of an invertible matrix $B$ are the new basis vectors
expressed in the old coordinates, then a point with new coordinates
$x_{\text{new}}$ has old coordinates
\[
  x_{\text{old}} = B\,x_{\text{new}},
  \qquad
  x_{\text{new}} = B^{-1}\,x_{\text{old}};
\]
for an orthonormal basis $B^{-1} = B^\top$, and the map is a rigid
transformation. Used in Chapter~\ref{ch:spectral_geometry} (eigenbases of
the Laplacian), Chapter~\ref{ch:geometry_data_ml} (PCA), and
Chapter~\ref{ch:surface_geometry} (local tangent frames).

\medskip\noindent\textbf{Orthonormal frame from a normal.}
\index{frame}
Given a unit normal $n$ in $\mathbb{R}^3$, an orthonormal frame
$(e_1, e_2, n)$ is obtained by taking any vector $a$ not parallel to $n$,
setting $e_1 = (a - \langle a, n\rangle n)/\|a - \langle a, n\rangle n\|$
and $e_2 = n \times e_1$. Used in Chapter~\ref{ch:curves_surfaces} and
Chapter~\ref{ch:applications_graphics} (tangent-space shading).

\index{coordinate transformations}\index{homogeneous coordinates}

\section{Barycentric Coordinates}
\label{sec:app_formulas_barycentric}

\medskip\noindent\textbf{Triangle barycentric coordinates.}
\index{barycentric coordinates}
With respect to an affine frame $A, B, C$ of $\mathbb{R}^2$, every point
$x$ has a unique representation $x = \lambda_A A + \lambda_B B + \lambda_C C$
with $\lambda_A + \lambda_B + \lambda_C = 1$; the coefficients solve
\[
  \begin{pmatrix} 1 & 1 & 1 \\ A_x & B_x & C_x \\ A_y & B_y & C_y \end{pmatrix}
  \begin{pmatrix} \lambda_A \\ \lambda_B \\ \lambda_C \end{pmatrix}
  = \begin{pmatrix} 1 \\ x_x \\ x_y \end{pmatrix}
\]
and are the signed area ratios
\[
  \lambda_A = \frac{A(\triangle xBC)}{A(\triangle ABC)},
  \quad
  \lambda_B = \frac{A(\triangle xCA)}{A(\triangle ABC)},
  \quad
  \lambda_C = \frac{A(\triangle xAB)}{A(\triangle ABC)},
\]
with the signed areas of Section~\ref{sec:app_formulas_areas_volumes}
(Section~\ref{sec:math_barycentric}). The point $x$ lies inside the
triangle iff all three coordinates are nonnegative, on the boundary iff
one is zero, and outside iff some coordinate is negative; this yields a
point-in-triangle test built from three orientation evaluations. Used in
Chapter~\ref{ch:math_foundations}, Chapter~\ref{ch:point_location}, and
Chapter~\ref{ch:sampling}.

\medskip\noindent\textbf{Linear interpolation.}
\index{interpolation}
Given scalar data $f_A, f_B, f_C$ at the vertices, the affine interpolation
at $x$ is
\[
  f(x) = \lambda_A f_A + \lambda_B f_B + \lambda_C f_C,
\]
the unique affine function agreeing with the data at the vertices; it is
affine-invariant and is the linear finite-element basis function on the
triangle (Section~\ref{sec:app_math_affine}). Used in
Chapter~\ref{ch:mesh_analysis} (finite elements, Section~\ref{sec:scientific_fem}),
Chapter~\ref{ch:applications_graphics} (texture and color interpolation),
and Chapter~\ref{ch:geometry_data_ml}.

\medskip\noindent\textbf{Tetrahedron barycentric coordinates.}
\index{barycentric coordinates}
For a tetrahedron $p_0 p_1 p_2 p_3$ the barycentric coordinates are the
signed volume ratios
\[
  \lambda_i = \frac{V(x\, p_j p_k p_l)}{V(p_0 p_1 p_2 p_3)},
  \qquad \{i, j, k, l\} = \{0, 1, 2, 3\},
\]
with $V(\cdot)$ the signed tetrahedron volume of
Section~\ref{sec:app_formulas_areas_volumes}; again $x$ is inside iff all
$\lambda_i \ge 0$, and interpolation $f(x) = \sum_i \lambda_i f_i$ is
affine. Used in Chapter~\ref{ch:mesh_analysis} (tetrahedral finite
elements) and Chapter~\ref{ch:mesh_generation}.

\medskip\noindent\textbf{Barycentric coordinates on the $d$-simplex.}
For an affine frame $p_0, \dots, p_d$ of $\mathbb{R}^d$ the coordinates of
$x$ are the unique solution of the $(d+1)\times(d+1)$ system
\[
  \sum_{i=0}^{d} \lambda_i = 1,
  \qquad
  \sum_{i=0}^{d} \lambda_i p_i = x,
\]
equivalently ratios of signed $d$-volumes, and the interpolation formula
$f(x) = \sum_i \lambda_i f_i$ is the general linear basis on simplices.
Generalized (non-affine, mean-value and harmonic) coordinates for arbitrary
polygons are used in the parameterization literature
(Chapter~\ref{ch:surface_modeling}).

\medskip\noindent\textbf{Face counts of a simplex.}
\index{simplex}\index{k-face}
The $d$-simplex has
\[
  f_k(\Delta^d) = \binom{d+1}{k+1}
\]
faces of dimension $k$ ($0 \le k \le d$), and
$\sum_{k=0}^{d} f_k = 2^{d+1} - 1$ proper faces; e.g., the tetrahedron has
$f_0 = 4$ vertices, $f_1 = 6$ edges, $f_2 = 4$ triangles, and $f_3 = 1$
body. Used in Chapter~\ref{ch:computational_topology} and
Chapter~\ref{ch:convex_polytopes} (Section~\ref{sec:polytopes_faces}).

\medskip\noindent\textbf{Euler--Poincar\'e formula.}
\index{Euler characteristic}\index{Euler--Poincar\'e formula}
For a convex $d$-polytope with $f_k$ faces of dimension $k$,
\[
  \sum_{k=0}^{d-1} (-1)^k f_k = 1 - (-1)^d,
\]
which for $d = 3$ is Euler's polyhedron formula $V - E + F = 2$ and for a
closed surface triangulated with Euler characteristic $\chi$ reads
\[
  V - E + F = \chi,
  \qquad
  \chi = 2 - 2g \text{ for an orientable closed surface of genus } g
\]
(Section~\ref{sec:math_topology}, Chapter~\ref{ch:computational_topology};
the planar $V - E + F = 2$ appears in Section~\ref{sec:appform_counting}).
The alternating sum is the discrete Euler characteristic of a simplicial
complex (Appendix~\ref{app:math_reference}, Section~\ref{sec:app_math_graphs}).
Used in Chapter~\ref{ch:half_edge} (Section~\ref{sec:halfedge_complexity}),
Chapter~\ref{ch:mesh_analysis}, and Chapter~\ref{ch:convex_polytopes}.

\begin{theorem}[Euler--Poincar\'e for Polyhedra]\label{thm:appform_euler_poincare}
  For a closed (boundaryless) polyhedral surface, $V - E + F$ is a
  topological invariant, equal to the Euler characteristic $\chi$; for an
  orientable surface triangulated with $n$ triangles,
  $E = \frac{3}{2}F$ and $V - \frac{1}{2}F = \chi$.
\end{theorem}

\index{barycentric coordinates}\index{Euler formula}
